\subsection{Which ordinary points give classical spectra?}\label{spreal:section}

Write \(k=\overline{\mathbb Q}_\ell\) and
\(\varepsilon=\varepsilon_\ell\). A representation is called strongly
irreducible if its restriction to every open subgroup is absolutely
irreducible. We work with continuous representations at actual
characteristic-zero points. They and their invariant lines can be defined
over a finite coefficient extension, by Lemma~\ref{newgal:coeff} and
extension of the coefficients of a finite basis.

\begin{lemma}[Finite self-twists at an ordinary place]\label{spreal:twists}
Let \(V\) be an irreducible representation of \(G_F\) of dimension \(n\).
Suppose the semisimplification on inertia at one place \(v\mid\ell\)
is a sum of characters \(\lambda_i\) whose restrictions to an open
subgroup of inertia are \(\varepsilon^{m_i}\) for distinct integers \(m_i\).
Every finite-order self-twist of \(V\) is unramified at \(v\).
If \(\varepsilon(I_v)=\varepsilon(G_F)\), then
\(V|_{G_{F(\mu_{\ell^\infty})}}\) is irreducible.
\end{lemma}
\begin{proof}
If \(V\simeq V\otimes\nu\), then multiplication by \(\nu|_{I_v}\)
permutes the \(\lambda_i\). On an open subgroup killing \(\nu\) the
characters remain distinct, since the cyclotomic image there is
infinite. The permutation is the identity, so \(\nu|_{I_v}=1\).
Put \(H=\ker\varepsilon\). The restriction to \(H\) is semisimple:
the sum of its simple submodules is nonzero and \(G_F\)-stable.
If it were reducible, its commutant would contain a nonzero traceless
endomorphism, obtained by subtracting a scalar from an equivariant
projection. The abelian quotient \(G_F/H\) acts on the space of such
endomorphisms by commuting operators. A common eigenvector \(T\ne0\)
gives an intertwiner with a character twist. Its kernel is
\(G_F\)-stable, so \(T\) is invertible. Schur's lemma and
\(\operatorname{tr}T=0\) exclude the trivial character. Taking
determinants shows that the character has order dividing \(n\).
It is continuous from its action on the line \(kT\). Triviality on
\(I_v\) and surjectivity onto the cyclotomic quotient give a contradiction.
\end{proof}

For the rank-three datum let \(K/\mathbb Q\) be the imaginary quadratic
field, \(\eta\) its character, and \(\Psi=\psi\varepsilon^e\).
Use the prime and character hypotheses in Section~\ref{rk3:sec}.
In particular \(\ell\nmid D_K h_K\operatorname{ord}(\psi)\),
\(\psi_\ell^2=1\), and \(\ell-1\nmid2e\). Thus \(\ell\ge5\).
The normalized matrix at complex conjugation is the identity with
negative sign. On \(G_K\) the residual matrices are diagonal with
characters \(\bar\Psi,1,\bar\Psi^{-1}\). Away from \(\ell\) the
prescribed inertia kernel kills elements with trivial residual action.
At \(\ell\) put \(J=I_\ell\cap\ker\psi_\ell\); here \(I_\ell\subset G_K\).
The ordinary equations of Section~\ref{ct:section} apply on \(J\).
\src{src/Development/AutomorphicArithmetic/AutomorphicVarieties/UnitaryThreeGroups/RhoKApply.lean}{76}{Defining equations for a characteristic-zero point}.

\begin{lemma}[The inertia flag]\label{spreal:flag}
On \(J\) the ordinary equations give an upper triangular flag with
ordered diagonal \((\varepsilon^e,1,\varepsilon^{-e})\). This holds
over \(k\) and over \(k[\delta]/(\delta^2)\) at deformations of such
a point. The flag is stable under every positive-sign element of the
decomposition group. Its graded inertia characters on \(I_\ell\)
are the indicated powers times characters of \(I_\ell/J\).
\end{lemma}
\begin{proof}
Choose \(s\in J\) with infinite cyclotomic value. The prescribed
characteristic polynomial has three distinct constant roots at \(s\).
Its Lagrange projectors split the free module into three free lines,
also over the dual numbers. In the ordered product of three shifted
matrices, put \(g\in J\) in position \(a\) and \(s\) in the other
two. In the eigenbasis the \((i,j)\) entry for \(i\ge a\ge j\)
is the corresponding entry of the \(a\)-th shifted matrix times
a unit. Taking \(a=j\) annihilates all lower entries and imposes
the stated diagonal. The common \(\varepsilon^e\)-eigenline is
uniquely specified by its eigenvalue at \(s\). The same description
in the quotient gives the next line. As \(J\) is normal and these
characters are conjugation invariant, the flag is preserved.
The extra inertia factors on the graded lines are trivial on \(J\).
\end{proof}

\begin{proposition}[The endoscopic block]\label{spreal:end}
Suppose the normalized rank-three point is of type II: on \(G_K\) it
preserves the middle coordinate and the outer coordinate plane,
and its action on the latter is irreducible. There is an odd
representation \(W\) of \(G_{\mathbb Q}\), with determinant \(\eta\),
such that the point restricts to \(1\oplus W|_{G_K}\). This \(W\)
is the normalized representation of a classical cuspidal newform
of weight \(2e+1\). Its automorphic representation has invariants
at principal level \(N|D_K|\).
\end{proposition}
\begin{proof}
On \(G_K\) denote by \(m\) the middle character and by \(d\) the
determinant of the whole representation. Both reduce to one. On \(J\)
the determinant is one by Lemma~\ref{spreal:flag}, and \(m\) is one
of its three constituents. The reduction of \(\varepsilon^e\) is
nontrivial on \(J\): otherwise the image of
\(\bar\varepsilon^e(I_\ell)\) has order at most two, contrary to
\(\ell-1\nmid2e\). Here \(\bar\varepsilon(I_\ell)=\mathbb F_\ell^\times\).
It follows that \(m|_J=1\). Since \([I_\ell:J]\le2\), reduction
detects both characters on \(I_\ell\). The argument holds at every
place of \(K\) over \(\ell\): conjugating from the marked place
either conjugates the matrices or takes their conjugate inverse
transposes, reversing and negating the three exponents together.
At primes away from \(\ell\), the inertia action on \(G_K\) factors
through its residual image by the kernel condition. This image has
order dividing \(\operatorname{ord}(\psi)\), prime to \(\ell\).
Thus \(m,d\) are unramified globally, and their images in the group
of principal units factor through \(\operatorname{Cl}(K)\). Since
\(\ell\nmid h_K\), both are trivial.

Let \(B(h)\in\mathrm{SL}_2(k)\) be the outer block for \(h\in G_K\).
Take \(W(h)=B(h)\) and \(W(c)=D=\operatorname{diag}(-1,1)\) at
complex conjugation. With \(J_2\) the antidiagonal unit matrix,
\(J_2B^{-t}J_2=DBD^{-1}\) whenever \(\det B=1\).
This is the action of \(c\) on the corner, and \(c^2=1\), giving the
claimed extension, determinant and oddness. It is continuous on the
two cosets and irreducible already on \(G_K\).

On \(J\) the highest inertia line of Lemma~\ref{spreal:flag} belongs
to the outer plane, since the middle character is one. In \(W\) it
is the unique common \(\varepsilon^e\)-eigenline, so normality of
\(J\) makes it stable under the full decomposition group. The two
graded characters on inertia have powers \(e,-e\) with finite parts
trivial on \(J\). Over the finite local extension killing the finite
parts the two characters are crystalline. Their ratio \(U\), upper
line divided by quotient, is an unramified twist of \(\varepsilon^{2e}\).
Over that local field \(L\) both \(H^0(L,U)\) and \(H^0(L,U^*(1))\)
vanish since \(2e\ge2\). The local Euler characteristic formula gives
\(\dim H^1(L,U)=[L:\mathbb Q_\ell]\) over the coefficient field.
The Bloch--Kato dimension formula gives the same dimension for
\(H^1_f(L,U)\), since \(\operatorname{Fil}^0 D_{\rm dR}(U)=0\).
The extension is therefore crystalline over \(L\) and de Rham over
\(\mathbb Q_\ell\), with distinct jumps \(\{-e,e\}\).

The representation is unramified outside \(\ell N D_K\).
Pan's regular Fontaine--Mazur theorem now gives a cuspidal eigenform
up to twist. Its weight is one plus the difference of the two jumps.
In the geometric Frobenius convention of Section~\ref{rk3:sec}, the
cohomological representation is \(W\varepsilon^{-e}\), with jumps
\(0,2e\). A finite character twist is absorbed into the newform.
Pan's only extra regular-case restriction for odd primes concerns
\(\ell=3\), already excluded here.
See \cite[Theorem 1.0.4]{panFM}.

For \(v\nmid\ell\), \(I_v\cap G_{K(\mu_N)}\) acts trivially by the
residual kernel condition. Thus inertia acts through a subgroup of
the finite abelian group \(\operatorname{Gal}(K(\mu_N)/\mathbb Q)\),
and its characters extend to that group. At \(\ell\), the finite
inertia characters are trivial on \(J\), hence are either trivial
or the prescribed quadratic local component, which likewise extends
with modulus dividing \(N\). Since \(K\subset\mathbb Q(\mu_{|D_K|})\),
Proposition~\ref{newgal:level} gives the claimed level. In particular
the classical trace and its nonvanishing comparison in
Lemmas~\ref{rk3:purity} and~\ref{rk3:integer} apply to every such block.
\end{proof}

\begin{proposition}[Strong irreducibility at the rank-three points]\label{spreal:strong}
The representation on \(G_K\) at every type I point and the outer
two-dimensional constituent at every type II point are strongly
irreducible. Both remain irreducible on \(G_{K(\mu_{\ell^\infty})}\).
\end{proposition}
\begin{proof}
Call the relevant irreducible representation \(V\) and put
\(d=\dim V\in\{2,3\}\). The inertia characters of
Lemma~\ref{spreal:twists} occur by Lemma~\ref{spreal:flag} and
Proposition~\ref{spreal:end}. As \(K\) is unramified at \(\ell\),
its local inertia already surjects onto \(\mathbb Z_\ell^\times\),
by total ramification of the local cyclotomic tower. This proves
the second assertion.

Suppose restriction to an open subgroup is reducible, and replace
that subgroup by its normal core \(H\). As in Lemma~\ref{spreal:twists},
the restriction to \(H\) is semisimple. Conjugation permutes its
isotypic spaces transitively, so their number times their common
multiplicity times the dimension of a simple constituent is \(d\).
Primality of \(d\) leaves two possibilities. A scalar restriction
to \(H\) contradicts the distinct characters on \(H\cap J\).
Otherwise there are \(d\) distinct lines permuted transitively.
Choose a stable lattice and replace it by the direct sum of its
intersections with these lines, enlarging coefficients if necessary.
In generators of these rank-one lattices the matrices are monomial
with unit coefficients. Their reduction has order prime to \(\ell\):
its diagonal subgroup lies in a product of finite residue-field
multiplicative groups and its permutation quotient has order
dividing \(d!\). Thus the reduction is semisimple. Independence of
residual semisimplification from the lattice now identifies it with
a sum of the given global characters. Its image is abelian, and
its permutation quotient, unchanged by reduction of the units, is
transitive abelian of prime degree. It is cyclic of order \(d\).

Let \(\nu\) be a faithful character of that quotient. The operator
which multiplies the successive lines in cyclic order by successive
powers of a primitive \(d\)-th root intertwines \(V\) with the
corresponding twist (up to replacing \(\nu\) by its inverse).
Lemma~\ref{spreal:twists} gives \(\nu|_{I_\ell}=1\). Its reduction
has order \(d\), and twisting must permute the residual constituents.
In dimension three those constituents are
\(1,\bar\Psi,\bar\Psi^{-1}\), so the orbit of \(1\) consists of all
three powers of \(\bar\nu\). In dimension two the two constituents
are \(\bar\Psi,\bar\Psi^{-1}\), so \(\bar\nu=\bar\Psi^2\).
In either case \(\bar\Psi^2|_{I_\ell}=1\), contrary to
\(\bar\Psi^2|_{I_\ell}=\bar\varepsilon^{2e}|_{I_\ell}\ne1\).
\end{proof}

\begin{corollary}[The polarized adjoint at an endoscopic point]\label{spreal:endrig}
For the characteristic-zero coefficient space \(A=M_3(k)\) with
the polarized adjoint action of a type II point,
\(H^1_f(\mathbb Q,A)=0\).
\end{corollary}
\begin{proof}
Decompose matrices into the outer endomorphism block, the middle
scalar, and the two off-diagonal rectangles taken together. All
three spaces are stable, since \(c\) acts by negative reversed
transpose. The outer block is the polarized adjoint of \(W|_{G_K}\).
Indeed its polarization matrix for this action is \(J_2\), the
product of the alternating determinant pairing
\(\left(\begin{smallmatrix}0&1\\-1&0\end{smallmatrix}\right)\)
and \(W(c)^{-1}\). Cyclic base change of the modular form to \(K\)
is cuspidal: a noncuspidal quadratic base change is induced from a
character of \(K\), which would give reducible restriction to
\(G_K\). The normalized Tate twist does not affect the adjoint.
Propositions~\ref{spreal:strong} and~\ref{newgal:adselmer} give its
Selmer vanishing. For the other summands restriction to \(K\) is
injective on rational cohomology by restriction and corestriction
and takes finite classes to finite classes. Over \(K\) the scalar
is trivial. A finite local class in \(H^1(K_v,k)\) at \(v\mid\ell\)
is unramified: the unramified line is crystalline and exhausts the
finite subspace by its dimension. A global everywhere unramified
additive character of \(G_K\) vanishes by class field theory.

The off-diagonal part over \(K\) is \(W|_{G_K}\oplus W^*|_{G_K}\).
These two representations are isomorphic since \(\det W|_{G_K}=1\).
For the weight \(k_0=2e+1\) newform just obtained, denote its
cohomological representation by \(V_f\). In Kato's convention
\(W=V_f(e)\). Here \(1\le e\le k_0-1\) and \(2e\ne k_0\).
Kato's Theorem~14.2(1) over the finite abelian extension \(K/\mathbb Q\)
gives finiteness of the discrete Bloch--Kato group for any stable
lattice in this twist. To pass to rational coefficients, the image of
the vector-space Selmer group under reduction in the discrete group
is divisible and hence zero. Its kernel is contained in the image
of integral \(H^1(G_{K,S},T)\), a finitely generated lattice module
for a finite \(S\) containing the ramification and \(\ell\).
A vector space contained in such an image is zero. This argument
may be made over the original finite coefficient field and then
extended to \(k\) as in Proposition~\ref{newgal:adselmer}.
See \cite[\S14.1 and Theorem 14.2(1)]{katoZeta}.
\end{proof}

The fixed ordinary equations place every continuous tangent cocycle
at this point in \(H^1_f\).
At a prime away from \(\ell\) its restriction to inertia factors
through a finite group, both for the point and the lift, by the
fixed kernel condition. Finite-group cohomology over \(k\) vanishes
in positive degree. At \(\ell\) pass to a finite local extension
inside the positive-sign subgroup killing the finite graded inertia
characters. By Lemma~\ref{spreal:flag} the lift over the dual numbers
has a stable ordered flag. Each rank-one graded module, regarded as
a two-dimensional representation over the coefficient field, is an
unramified representation tensored with its fixed cyclotomic power
and is crystalline. Successive extensions, taken in the order of the
flag, are crystalline as well. Indeed the Hom from a later graded
piece to the already constructed part has a filtration with factors
unramified characters times \(\varepsilon^s\) for \(s\ge e\ge1\).
It is crystalline and its de Rham filtration in degree zero vanishes.
When \(e\ge2\), distinct exponents alone imply that its dual twist
has no invariants. For \(e=1\) we use the modular block at the
coefficient prime. Its normalized Weil--Deligne parameter has zero
monodromy by potential crystallinity in Proposition~\ref{spreal:end}.
The full comparison at \(\ell\) in Proposition~\ref{newgal:level}
identifies it in the weight-zero normalization with the local
parameter of the unitary newform. This parameter is tempered by
Blasius's Theorem~1 for holomorphic forms of weight at least two.
With zero monodromy its Weil Frobenius eigenvalues have magnitude
one, also after finite local restriction. The middle character is one.
Every ratio of two graded characters after setting \(\delta=0\)
thus has Weil Frobenius eigenvalue of magnitude one, whereas
\(\varepsilon\) gives \(q_L^{-1}\) on geometric Frobenius.
No such ratio is \(\varepsilon\), so the dual twist has no invariants
in this case either.
See \cite[Theorem 1]{blasiusRam}.
The same local dimension
calculation as in Proposition~\ref{spreal:end} gives \(H_f^1=H^1\).
The lifted representation and its reduction being crystalline there,
their extension class becomes a coboundary over \(B_{\rm cris}\).
Restriction on cohomology with these period coefficients is injective
for finite local extensions by corestriction. The global tangent
class is therefore finite. Corollary~\ref{spreal:endrig} makes it a
coboundary, giving infinitesimal rigidity modulo conjugation at
these endoscopic points.

\begin{proposition}[Finite tangent classes on the rank-four component]\label{spreal:four}
For the good rank-four datum, at every actual characteristic-zero
point every tangent arising from the specified component gives a
finite local extension class in the adjoint at every finite prime.
\end{proposition}
\begin{proof}
Put \((\theta_i)=(\alpha\Psi,\Psi,\alpha,1)\). The ordinary
equations in Proposition~\ref{dp:local} give a flag with graded inertia
characters \(\theta_i\), also over the dual numbers. The proof with
projectors there applies at any point: the distinguished inertia
matrix has the same four distinct eigenvalues \(d_i=\theta_i(s_0)\).
After killing the finite inertia factors the four powers of
\(\varepsilon\) are \(e+1,e,1,0\). Denote the characteristic-zero
graded characters on the full decomposition group by \(\kappa_i\).

Let \(\sigma=s_\ell\) be the specified Frobenius lift with
\(\varepsilon(\sigma)=1\). We claim
\(\overline{\kappa_i(\sigma)}=\overline{\theta_i(\sigma)}\)
with the indicated labels, even if some residual eigenvalues coincide.
The component hypothesis provides a complete local domain \(B\)
with an ordinary matrix family \(R\), specializing to the diagonal
point and to the given point. Invert the nonzero constant
\(a=\prod_{i<j}(d_i-d_j)\). On \(B[1/a]\) the projectors
\(P_i\) of \(R(s_0)\) have rank one. Indeed their ranks are locally
constant on this connected spectrum and are one at the diagonal
specialization. The triangular equations make
\(b_i=\operatorname{tr}(P_iR(\sigma))\) the scalar on the
corresponding flag quotient, so that \(b_i\) is a root of the
characteristic polynomial of \(R(\sigma)\).
Over the Henselian ring \(B\) that polynomial splits into monic
factors \(F_u\), one for each distinct residual eigenvalue \(u\),
with \(\bar F_u=(X-u)^{n_u}\). In the domain \(B[1/a]\)
exactly one \(F_u(b_i)\) vanishes, by coprimality. The diagonal
specialization shows that its label is
\(u=\overline{\theta_i(\sigma)}\). At the other specialization
the same factor vanishes, proving the claim. In this argument any
\(O\)-algebra specialization \(B\to\overline{\mathbb Q}_\ell\)
is integral and reduces to the given residual map. Indeed an
element of negative valuation has a polynomial relation over
\(O\) with constant term one and all higher coefficients in
\(\mathfrak m_O\), by using the minimal polynomial of its inverse
over \(E\). Such a relation gives a unit in \(B\). For an element
of \(\mathfrak m_B\) mapping to a unit, its image has monic
minimal polynomial with unit constant term, again impossible.

Let \(I_0=I_\ell\cap\ker\psi\cap\ker\omega\). In the quotient of
the local group by \(I_0\) take the closed subgroup generated by
\(\sigma\) and call its preimage \(G_L\). The generated procyclic
subgroup maps isomorphically to \(\widehat{\mathbb Z}\), since its
generator maps to a Frobenius generator there. Thus \(L\) is finite,
\(G_L\cap I_\ell=I_0\), and \(\sigma\in G_L\).
For any higher and lower graded characters, their ratio differs from
\(\varepsilon\) on this group. If their exponents do not differ
by one this follows on \(I_0\). The ratios with exponent difference
one are those with underlying diagonal ratio \(\alpha\), or
\(\Psi/\alpha\) if \(e=2\). Since \(\bar\alpha(\sigma)=
\bar\xi(\sigma)\), the good-datum condition
\(\bar\xi(\sigma)\bar\Psi(\sigma)^i\ne1\) for \(i=0,-1\)
proves the assertion by the labeled congruence.

The dual-number flag reduces modulo \(\delta\) to this same flag,
by the distinct labeled eigenvalues at \(s_0\). On \(G_L\) its
graded pieces are unramified two-dimensional representations times
their cyclotomic powers. For successive extensions from a later
piece to the already constructed part, the Hom has strictly positive
cyclotomic exponents and its dual twist has no invariants by the
preceding test. As in the rank-three calculation the local Euler
and Bloch--Kato formulas give \(H^1=H^1_f\), inductively giving
crystallinity over \(L\) and finiteness of the extension class over
\(\mathbb Q_\ell\). At every other prime the point and its tangent
kill a common finite-index subgroup of inertia by the prescribed
kernel conditions. Averaging over the finite quotient gives the
unramified cohomology condition.
\src{src/Development/AutomorphicArithmetic/AutomorphicVarieties/SymplecticGSp4/DecompSigmaOrd.lean}{27}{The domain and the two specializations},
\src{src/Development/AutomorphicArithmetic/AutomorphicVarieties/PadicHodgeTheory/SplAdjUnit.lean}{264}{the labeled Frobenius congruence}
and \src{src/Development/AutomorphicArithmetic/AutomorphicVarieties/SymplecticGSp4/TbTriDiag.lean}{310}{the local nonresonance test}.
\end{proof}

For each irreducible point choose an imaginary quadratic \(K\)
unramified at \(\ell\) with irreducible restriction as in
Lemma~\ref{rk4:quadratic}. Its local inertia has full cyclotomic
image. Thus Lemma~\ref{spreal:twists} verifies the cyclotomic
irreducibility needed by Proposition~\ref{newgal:adselmer}.
A regular polarized cuspidal attachment over this \(K\) therefore
suffices for rank-four rigidity by Proposition~\ref{spreal:four}
and the symplectic projection of
Section~\ref{newgal:section}.
\src{src/Development/AutomorphicArithmetic/AutomorphicVarieties/SymplecticGSp4/MapRepApply.lean}{99}{Polynomial projectors in the ordinary local deformation condition}.

\begin{lemma}[The excluded rank-four block]\label{spreal:even}
Suppose a rank-four point preserves the two coordinate planes of
the block condition. If its representation \(V\) on coordinates
\(1,3\) (one-based indices) is irreducible, it is strongly
irreducible, potentially crystalline with distinct weights at
\(\ell\), and even. Its residual semisimplification is
\(\bar\alpha\bar\Psi\oplus\bar\alpha\).
\end{lemma}
\begin{proof}
The integral specialization gives the asserted residual characters.
At complex conjugation their common value is \(-1\), since
\(\Psi(c)=1\), \(\xi(c)=-1\) and \(\bar\omega=\bar\varepsilon\).
An involution reducing to \(-I\) is itself \(-I\) in odd residual
characteristic, by the idempotent \((1+c)/2\). The block inherits
potential crystallinity from Proposition~\ref{spreal:four}'s
local argument. Intersecting its plane with the full flag gives
two distinct exponents on open inertia, since all four are distinct.
The image of inertia is finite away from \(\ell\).

If \(V\) becomes reducible on an open subgroup, pass to its normal core.
The scalar case is excluded by the exponents. The other possibility,
by the isotypic decomposition in Proposition~\ref{spreal:strong},
is induction of a character \(\lambda\) from a quadratic field \(L\).
This field is real since \(c\) fixes both lines. If \(\ell\) is
nonsplit in \(L\), its decomposition group exchanges the two lines.
Their characters on a common open inertia subgroup then have the
same cyclotomic exponent, as conjugation preserves the cyclotomic
character. If \(\ell\) splits, write \(m_1,m_2\) for the exponents
of \(\lambda\) at the two places. Take a nontrivial totally positive
unit \(u\) of norm one in \(L\). Replacing it by a power kills its
image under the finite inertia characters at every other ramified
place, and puts it in the open unit subgroups with the specified
cyclotomic powers at the two places over \(\ell\). Global reciprocity
on the principal idele of \(u\) gives
\(\iota_1(u)^{m_1}\iota_2(u)^{m_2}=1\), up to simultaneous
inversion according to the reciprocity convention. As
\(\iota_1(u)\iota_2(u)=1\) and \(u\) has infinite order, this gives
\(m_1=m_2\), again impossible. This proves strong irreducibility.
\end{proof}

The rank-four block exclusion would therefore follow from the even
two-dimensional Fontaine--Mazur assertion for these ordinary, residually
split representations. The oddness in Pan's theorem fails precisely
on this block. Calegari's ordinary theorem assumes
\(\ell>7\) and residual image containing \(\mathrm{SL}_2(\mathbb F_\ell)\).
The displayed residual semisimplification rules out this image for
every choice of lattice. See \cite[Theorem 1.2]{calegariEven}.

The missing higher-rank assertion is cuspidal attachment for every
relevant characteristic-zero irreducible point. For example, the
residually reducible lifting theorem of Allen, Newton and Thorne
requires in its Theorem~1.1(4) that the distinct residual constituents
be individually conjugate self-dual with the same multiplier.
After a common twist the actual rank-three constituents have ratios
\(\bar\Psi,\bar\Psi^{-1}\) to the middle character, both pulled back
from \(\mathbb Q\). Individual polarization would make each ratio
self-inverse, thereby identifying the two outer characters, contrary
to the distinctness hypothesis. This obstruction persists under
restriction to a CM extension and any common character twist:
either the square remains nontrivial or the outer characters coincide.
In rank four the residual characters up to twist are
\(1,a,b,ab\), with \(a=\bar\alpha\) and \(b=\bar\Psi\), again
restricted from \(\mathbb Q\). If condition (4) holds, \(a,b\)
are independent quadratic characters. For \(H=\ker a\) the sum
of these four characters is \(\operatorname{Ind}_{H}^{G_F}(1\oplus b|_H)\),
contradicting the residual primitivity in (7). This is also invariant
under common twisting and applies after any CM restriction.
Conditions (5)--(6) of the same theorem independently require a place
away from \(\ell\) with consecutive unramified cyclotomic twists in
the characteristic-zero local semisimplification, and an ordinary
congruent cuspidal seed which is a Steinberg twist there.
See \cite[Theorem 1.1]{antLifting}.

The precise rank-three statement used by the source gives a conjugate
self-dual cuspidal attachment of \(r|_{G_K}\varepsilon^{-1}\) for
every ordinary point over a finite coefficient DVR whose restriction
to every open subgroup of \(G_K\) is irreducible. Proposition
\ref{spreal:strong} supplies this last hypothesis. The proposed passage
uses an auxiliary CM frame, arbitrarily large weights in a prescribed
residue class supporting adjoints of rank-two seeds, propagation in the
ordinary determinant space to the specified point, and descent to
\(K\). To reach the given point one must show its restriction lies
on the automorphic quotient and then descend the attachment. Existence
of a seed at each high weight supplies only points on the seed locus.
\src{src/Development/AutomorphicArithmetic/AutomorphicVarieties/ModularityArgument/ThresholdArguments/DenseThresholdCases/AutEnd3AltAlt.lean}{20}{The ordinary-point assertion},
\src{src/Development/AutomorphicArithmetic/AutomorphicVarieties/ModularityArgument/ThresholdArguments/DenseThresholdCases/SymUse.lean}{27}{the cofinal-weight seed assertion}
and \src{src/Development/AutomorphicArithmetic/AutomorphicVarieties/ModularityArgument/ThresholdArguments/DenseThresholdCases/GdEndOf.lean}{22}{the linked-point and descent inputs}.

For rank four the asserted output gives, for every imaginary quadratic
\(K\) on which the restriction stays irreducible, the cuspidal
representation and the prescribed comparison witness. On a sufficiently
deep auxiliary tower the proposed proof constructs an ideal \(J\)
in a universal determinant algebra. Specializations killing \(J\)
at the prescribed weights must be pure and, in the irreducible
case, automorphic after descent. To use these implications at the
given point one also needs a specialization killing \(J\) which
realizes its restriction, up to conjugacy, with the labeled inertia
characters. In the case with four seed characters the chain assertion
upgrades membership of some component over a residual specialization
to membership of all components after the stated dimension bound
\(30(C+1)^2\le d\). Establishing the existence of \(J\) with those
arithmetic properties and the coverage of the specific point remains
necessary to convert the finite orbit estimate into a comparison for
the entire Selmer module.
\src{src/Development/AutomorphicArithmetic/AutomorphicVarieties/ModularityArgument/AutomorphyLifting/IsAutCsdA.lean}{16}{The quantifier over quadratic restrictions},
\src{src/Development/AutomorphicArithmetic/AutomorphicVarieties/ModularityArgument/ThresholdArguments/DenseThresholdCases/PotentialModularityTower.lean}{63}{the ideal and its required specialization},
\src{src/Development/AutomorphicArithmetic/AutomorphicVarieties/ModularityArgument/ThresholdArguments/DenseThresholdCases/PromodularityCaseThree.lean}{26}{the component propagation}
and \src{src/Development/AutomorphicArithmetic/AutomorphicVarieties/ModularityArgument/ThresholdArguments/DenseThresholdCases/ModularityRankFourOfPromodularity.lean}{20}{the implication to finiteness, block exclusion, and automorphy}.
